\noindent\resizebox{\textwidth}{!}{%
\begin{tikzpicture}[every node/.style={font=\footnotesize}]

% Three stage headers
\node[dia/header=4.5cm, minimum height=1.0cm] (s1h)
  {\small\textbf{Етап 1: Теоретико-аналітичний}};
\node[dia/header=4.5cm, minimum height=1.0cm, right=0.8cm of s1h] (s2h)
  {\small\textbf{Етап 2: Практико-розробницький}};
\node[dia/header=4.5cm, minimum height=1.0cm, right=0.8cm of s2h] (s3h)
  {\small\textbf{Етап 3: Інтегративно-рефлексивний}};

% Arrows between stages
\draw[dia/arrow] (s1h.east) -- (s2h.west);
\draw[dia/arrow] (s2h.east) -- (s3h.west);

% Stage 1 content boxes (placed first so labels can reference them)
\node[dia/lrbox=4.5cm, below=0.7cm of s1h.south, anchor=north] (s1r1) {%
  аналіз потреб,\\формулювання вимог};
\node[dia/lrbox=4.5cm, below=0.5cm of s1r1, anchor=north] (s1r2) {%
  когнітивне навантаження\\(scaffolding)};
\node[dia/lrbox=4.5cm, below=0.5cm of s1r2, anchor=north] (s1r3) {%
  TK + CK\\(базові знання)};
\node[dia/lrbox=4.5cm, below=0.5cm of s1r3, anchor=north] (s1r4) {%
  модулі 1--6\\(теорія WebAR, інструменти)};

% Stage 2 content boxes (aligned to Stage 1 rows)
\node[dia/lrbox=4.5cm, anchor=north] at (s2h.south |- s1r1.north) (s2r1) {%
  вибір ПЗ, планування,\\розробка};
\node[dia/lrbox=4.5cm, anchor=north] at (s2h.south |- s1r2.north) (s2r2) {%
  втілення, агентність,\\присутність};
\node[dia/lrbox=4.5cm, anchor=north] at (s2h.south |- s1r3.north) (s2r3) {%
  TPK, TCK\\(інтеграція)};
\node[dia/lrbox=4.5cm, anchor=north] at (s2h.south |- s1r4.north) (s2r4) {%
  модулі 7--12\\(трекінг, 3D, ML)};

% Stage 3 content boxes (aligned to Stage 1 rows)
\node[dia/lrbox=4.5cm, anchor=north] at (s3h.south |- s1r1.north) (s3r1) {%
  впровадження,\\оцінювання};
\node[dia/lrbox=4.5cm, anchor=north] at (s3h.south |- s1r2.north) (s3r2) {%
  саморегуляція,\\мотивація};
\node[dia/lrbox=4.5cm, anchor=north] at (s3h.south |- s1r3.north) (s3r3) {%
  TPACK\\(повна інтеграція)};
\node[dia/lrbox=4.5cm, anchor=north] at (s3h.south |- s1r4.north) (s3r4) {%
  модулі 13--18\\(проєкти, рефлексія)};

% Row labels on the left (positioned relative to Stage 1 content boxes)
\node[font=\footnotesize\bfseries\itshape, anchor=east] at ($(s1r1.west) + (-0.4cm,0)$) {Фази моделі ІОР:};
\node[font=\footnotesize\bfseries\itshape, anchor=east] at ($(s1r2.west) + (-0.4cm,0)$) {Фактори CAMIL:};
\node[font=\footnotesize\bfseries\itshape, anchor=east] at ($(s1r3.west) + (-0.4cm,0)$) {TPACK:};
\node[font=\footnotesize\bfseries\itshape, anchor=east] at ($(s1r4.west) + (-0.4cm,0)$) {Модулі курсу:};

% Dominant condition label at bottom
\node[dia/rbox=15.1cm, below=0.6cm of s2r4.south, font=\footnotesize] (cond) {%
  \textbf{Провідна умова:}~~психолого-педагогічна \textbar{} технологічна + методична \textbar{} організаційна%
};

\end{tikzpicture}%
}
